\section{A bound in terms of the endpoint}\label{sec:count}

Theorem~\ref{thm:family} takes a desired count $B$ and produces
an endpoint. We now reverse that description. For every even $s$,
it gives the computable bound
\begin{equation}\label{eq:inverse}
 I(s)\ge\max\bigl(\{0\}\cup
       \{B\ge1:D_B(66+45\log B)\le s\}\bigr).
\end{equation}
The ceiling disappears here because $s$ is even. To obtain the
simpler formula in Theorem~\ref{thm:main}, we need an explicit
estimate for $D_B$ and control of the intervals between endpoints.

\subsection{The totient estimate and the square-root constant}

\begin{lemma}\label{lem:totients}
For every integer $k\ge1$, with $a=3/\pi^2$,
\begin{equation}\label{eq:totient-bound}
 D_k\le ak^2+k(\log k+2).
\end{equation}
\end{lemma}
\begin{proof}
Counting coprime ordered pairs in $\{1,\ldots,k\}^2$ gives
\[
 2D_k-1=\sum_{d=1}^k\mu(d)\lfloor k/d\rfloor^2,
\]
where $\mu$ is the M\"obius function. The subtraction of $1$
accounts for the only coprime diagonal pair, $(1,1)$. Replacing
each squared floor by $(k/d)^2$ costs at most
$2k\sum_{d=1}^k1/d$. Extending the resulting M\"obius series
to infinity costs at most $k^2\sum_{d>k}d^{-2}\le k$.
Since $\sum_{d\ge1}\mu(d)/d^2=6/\pi^2$, we obtain
\[
 D_k\le ak^2+k\sum_{d=1}^k\frac1d+\frac k2+\frac12
       \le ak^2+k(\log k+2).
\]
\end{proof}

The same absolute-error estimate gives
$D_B=aB^2+O(B\log B)$.
The coefficient $a=3/\pi^2$ is half the density $6/\pi^2$ of
coprime pairs, since $D_B$ counts fractions with numerator no
greater than denominator. Thus the leading endpoint in
\eqref{eq:family} is
\[
 s_B\sim45aB^2\log B.
\]
The coefficient $45$ records the balance used in
Section~\ref{sec:endpoints}: the denominator cost is at most
$36D_B\log B$, while the decay gain per unit of $S=s+1$ is greater
than $4/5$. Their ratio gives $45=36/(4/5)$.

Since $\log B\sim\tfrac12\log s_B$, solving for $B$
gives the constant
\[
 \sqrt{\frac{2}{45a}}=\pi\sqrt{\frac2{135}}=c
\]
in front of $\sqrt{s/\log s}$. The factor $2$ comes from the
logarithm of the quadratic term $B^2$. We next show that this
leading constant already satisfies the required finite inequalities.

\subsection{The unrounded bound for large endpoints}

\begin{proposition}\label{prop:unrounded}
For every even integer $s\ge42\,170$,
\[
 I(s)\ge c\sqrt{\frac{s}{\log s}}.
\]
\end{proposition}
\begin{proof}
The endpoints $s_k$ increase strictly. Between $s_{k-1}$ and
$s_k$, the family guarantees at least $k-1$ irrational values.
Since $x/\log x$ increases for $x>\ep$, it suffices to check
the largest even point in each interval, namely $s_k-2$:
\begin{equation}\label{eq:unrounded-endpoint-test}
 c^2(s_k-2)<(k-1)^2\log(s_k-2),\qquad k\ge26.
\end{equation}
The first interval starts at $s_{25}=42\,170$.

We first prove \eqref{eq:unrounded-endpoint-test} for every
$k\ge256$. We record the error in the totient estimate as a
relative factor $q_k$ and use it to form an upper bound
$\mathcal U(k)$ for $s_k-2$. Put
\[
 \ell=\log k,\qquad b_0=\frac{66}{45}=\frac{22}{15},\qquad
 q_k=1+\frac{\ell+2}{ak},\qquad
 \mathcal U(k)=45ak^2(\ell+b_0)q_k.
\]
Lemma~\ref{lem:totients} gives $D_k\le ak^2q_k$. Together with
the ceiling inequality $2\lceil x/2\rceil-2<x$, this yields
\begin{equation}\label{eq:endpoint-upper}
 s_k-2<D_k(66+45\ell)\le\mathcal U(k).
\end{equation}
It is therefore enough to prove
$c^2\mathcal U(k)<(k-1)^2\log\mathcal U(k)$. Since
$c^2=2/(45a)$, this is
\begin{equation}\label{eq:upper-comparison}
 2(\ell+b_0)\frac{q_k}{(1-1/k)^2}<\log\mathcal U(k).
\end{equation}

Separate the common term $2(\ell+b_0)$ on both sides.
The excess on the left comes from the totient error and from
having only $k-1$ guaranteed values before the next endpoint.
We bound this cost by
\[
 \mathcal T(k)=
 \frac{2(\ell+b_0)((\ell+2)/a+2)}{k(1-1/k)^2},
\]
because
\[
 \frac{q_k}{(1-1/k)^2}-1
   =\frac{((\ell+2)/a+2)/k-1/k^2}{(1-1/k)^2}.
\]
The excess on the right is the margin available to absorb this cost:
\[
 \mathcal G(k)+\log q_k,\qquad
 \mathcal G(k)=\log(45a(\ell+b_0))-2b_0.
\]
As $q_k>1$, it suffices to show $\mathcal T(k)<\mathcal G(k)$.
The margin $\mathcal G$ increases. We will show that the error
bound $\mathcal T$ decreases for $k\ge256$, so a comparison at
$256$ will settle every $k\ge256$. Its logarithmic
derivative with respect to $\ell$ is
\[
 \frac1{\ell+b_0}+\frac1{\ell+2+2a}-1-\frac2{k-1}<0
                       \qquad(k\ge256).
\]
Indeed, $\ell>11/2$, and the sum of the first two terms is
less than $2/11+2/15<1$. It remains to compare the functions
at $k=256$.

Use $\pi<22/7$ and $11/2<\log256<111/20$. Direct rational
comparison gives
\begin{align}
 \mathcal T(256)
 &<2\left(\frac{111}{20}+\frac{22}{15}\right)
   \left(\left(\frac{111}{20}+2\right)\frac{484}{147}+2\right)
   \frac{256}{255^2}<\frac32,\label{eq:tail-T}\\
 45a(\log256+b_0)&>
       45\frac{147}{484}\left(\frac{11}{2}+\frac{22}{15}\right)>95.
       \label{eq:tail-G}
\end{align}
Moreover, $\ep<11/4$ and $(11/4)^9<95^2$, so
$\log95>9/2$. Thus
\[
 \mathcal G(256)>\frac92-\frac{44}{15}=\frac{47}{30}>\frac32.
\]
Monotonicity proves \eqref{eq:upper-comparison} throughout the
infinite range $k\ge256$.

For the remaining $230$ indices $26\le k\le255$, use the exact
totient sums in \eqref{eq:family}. The rational logarithm bounds
specified in Appendix~\ref{app:endpoint-checks} verify
\[
 \frac{(k-1)^2\,\underline{\log}(s_k-2)}{s_k-2}
       -\frac{968}{6615}>\frac{24529}{10^8}.
\]
Here $\underline{\log}(x)\le\log x$ is the explicit rational
lower bound, and $968/6615=2(22/7)^2/135>c^2$.
This proves \eqref{eq:unrounded-endpoint-test} for all remaining
indices, completing the proposition.
\end{proof}

\subsection{Keeping the floor down to the two-value count}

It remains to cover $160\le s<42\,170$. Let $t_k$ be the
sharper endpoint in Section~\ref{sec:endpoints} for $2\le k\le10$,
and let $t_k=s_k$ for $11\le k\le25$. The resulting increasing
list is
\begin{center}
\begin{tabular}{@{}rr@{\qquad}rr@{}}
\toprule
$k$&$t_k$&$k$&$t_k$\\
\midrule
2&160&14&11\,826\\
3&364&15&13\,528\\
4&588&16&15\,262\\
5&1030&17&18\,576\\
6&1296&18&20\,000\\
7&1992&19&23\,820\\
8&2502&20&25\,704\\
9&3260&21&28\,422\\
10&3828&22&30\,766\\
11&7306&23&35\,622\\
12&8180&24&37\,624\\
13&10\,524&25&42\,170\\
\bottomrule
\end{tabular}
\end{center}
On $t_{k-1}\le s<t_k$, we have $I(s)\ge k-1$. To prove the
rounded bound, it is enough for the expression inside the floor
to stay strictly below $k$. Once again, its largest value on the
even endpoints occurs at $s=t_k-2$. The required test is now
\begin{equation}\label{eq:rounded-endpoint-test}
 c^2(t_k-2)<k^2\log(t_k-2),\qquad 3\le k\le25.
\end{equation}
The $23$ exact comparisons in Appendix~\ref{app:endpoint-checks}
give the stronger bound
\[
 \frac{k^2\,\underline{\log}(t_k-2)}{t_k-2}
        -\frac{968}{6615}>\frac{14325}{10^8}.
\]
It follows that
$\lfloor c\sqrt{s/\log s}\rfloor\le k-1\le I(s)$ on every
one of these intervals. They cover all even integers from $160$
through $42\,168$, and Proposition~\ref{prop:unrounded} covers
the rest. This proves Theorem~\ref{thm:main}.

The floor also gives a completely explicit error form:
\[
 I(s)>\left(c-\sqrt{\frac{\log s}{s}}\right)
                \sqrt{\frac{s}{\log s}},\qquad s\ge160\text{ even}.
\]
The coefficient is positive throughout this range. In fact,
$c>19/50$ and $\sqrt{\log160/160}<9/50$, while $\log s/s$
decreases for $s>\ep$. Hence the coefficient is greater than $1/5$.
